% C.3bv. The complete NK proof passed root and both independent peer audits on 19 November 2026.
\paragraph{C.3bv. Whole proper $K_1$-family lifting (Proved).}
Let $k$ be algebraically closed, let $(R,\mathfrak m)$ be a noetherian
local coefficient-$k$ algebra with residue field $k$, and let $P/k$
be a smooth connected projective curve.  Write
$B=\operatorname{Spec}R\times_kP$.  Rank and determinant give an
integral isomorphism
\begin{equation}
 K_0(B)\xrightarrow[\sim]{\ (\operatorname{rk},\det)\ }
                   \mathbf Z\oplus\operatorname{Pic}(B).
                                                        \tag{NK.1}
\end{equation}
If $I$ is its rank-zero ideal, then
\begin{equation}
 I^2=F_\gamma^2K_0(B)=0,\qquad
       (\psi^q-1)(\psi^q-q)=0\quad(q\geq1).
                                                        \tag{NK.2}
\end{equation}
These assertions require no normality, regularity, henselian property,
or characteristic hypothesis on $R$.

Suppose further that $\operatorname{char}k=p>0$, $R$ is a normal
domain, and $g:Z\to\operatorname{Spec}R$ is a proper birational map
from a regular integral noetherian scheme, and an isomorphism away
from the closed point.  Put $E=(Z\times_Rk)_{\mathrm{red}}$.
Then restriction is surjective on the entire fixed-product group:
\begin{equation}
 K_1(Z\times_kP)[1/p]\longrightarrow K_1(E\times_kP)[1/p]
                         \quad\text{is surjective}.
                                                        \tag{NK.7}
\end{equation}
This conclusion is dimension-independent and requires neither
henselianity nor an ample exceptional conormal.  It concerns classes
in perfect $K$-theory, not extension of specified automorphisms as
objects.  The pointwise kernel and exact-endpoint consequence are
given in (NK.9a)--(NK.10) below.

\emph{1. Generation by line bundles on the actual proper product.}
Choose an ample line bundle $H$ on $P$ and use the same letter for
its pullback to $B$.  Let $V$ be a vector bundle of rank $r\geq2$.
For $n$ sufficiently large, relative Serre generation gives finitely
many global sections generating $V\otimes H^n$
\citep[Lemma 30.17.1]{frC3azStacksSerre}.
Their restrictions generate the bundle on the closed fibre $P$.
In their finite-dimensional coefficient space, vanishing at a
specified point imposes $r$ independent linear conditions.  The
incidence over the one-dimensional $P$ therefore has dimension at
most the dimension of the coefficient space plus $1-r$, strictly
smaller than that dimension.  Its image is closed because $P$ is
proper.  A general $k$-linear combination is consequently nowhere
zero on $P$.  The coefficient field $k\subset R$ lifts this
combination to an actual section on $B$.

Its zero scheme is closed and proper over the local scheme
$\operatorname{Spec}R$.  A nonempty image would be closed and
contain the closed point, contradicting the closed-fibre choice.
Thus the section is nowhere zero on all of $B$.  It gives a line
subbundle $H^{-n}\hookrightarrow V$ with locally free quotient.
Induction gives a filtration with line-bundle quotients.  Since $B$
has an ample line bundle over an affine base, perfect $K_0(B)$ is
the vector-bundle Grothendieck group
\citep[Corollary 3.9 and Proposition 3.10]{D_ThomasonTrobaugh1990}.
It follows that line-bundle classes generate $K_0(B)$ additively.
The argument does not assert that the filtration splits.

\emph{2. Finite Cartier supports and the square-zero ideal.}
Given a line bundle $M$ on $B$, take $n$ so large that
$M\otimes H^n$ and $H^n$ are generated.  Choose sections whose
restrictions to the integral closed curve are nonzero.  Their
closed-fibre zeros are finite effective Cartier divisors, possibly
empty.  Let $D_+$ and $D_-$ be the actual zero schemes on $B$.
They are locally principal.  The fibrewise Cartier criterion
\citep[Lemma 31.19.9(3)]{frNKStacksRelativeCartier} makes each
an effective Cartier divisor flat over $R$ in a neighbourhood of
every closed-fibre point.  At a point where the section is a unit,
the assertion is immediate.  These neighbourhoods cover $B$:
their closed complement is proper over $R$ and misses its closed
fibre.  Hence $D_\pm$ are relative Cartier divisors everywhere.

The divisors are proper over $R$ with finite closed fibre.
The proper finiteness-near-a-fibre lemma makes them finite over
an open neighbourhood of the closed point
\citep[Lemma 37.44.2, using Lemma 37.44.1]
{frNKStacksFiniteNearFibre,frNKStacksProperFinite}.
Such an open in a local scheme is the whole scheme.  Thus both
divisors are finite over $R$, and their rings are semilocal.
Every line bundle on a semilocal ring is trivial: choose residue-field
bases at its finitely many maximal ideals, lift them simultaneously
by the Chinese remainder theorem, and apply Nakayama's lemma at
each maximal ideal.  The resulting map from the ring to the
rank-one projective module is an isomorphism at every maximal ideal,
and hence an isomorphism
\citep[Lemma 10.78.7 and printed hints]{frWPStacksSemilocal}.

For a section of a line bundle $L$ with Cartier zero divisor $D$,
the Cartier sequence gives
$[L]-1=[L][\mathcal O_D]$.  Here $L|_D$ is trivial, so the right
side is $[\mathcal O_D]$.  Apply this to $L=M H^n,H^n$, subtract,
and multiply by $H^{-n}$.  The latter twist is also trivial on
both finite semilocal supports.  We obtain the exact identity
\begin{equation}
 [M]-1=[\mathcal O_{D_+}]-[\mathcal O_{D_-}]
                       \quad\text{in perfect }K_0(B).
                                                        \tag{NK.3}
\end{equation}
For any other line bundle $N$, its restrictions to $D_\pm$ are
trivial.  Perfect Cartier pushforward and the projection formula
therefore give
\begin{equation}
                 ([M]-1)([N]-1)=0.                    \tag{NK.4}
\end{equation}
Only Cartier divisor sheaves, which are perfect, were used; no
identification of arbitrary coherent and perfect classes occurs.

Line generation and (NK.4) show $I^2=0$.  Consequently
$[MN]-1=([M]-1)+([N]-1)$, so $M\mapsto[M]-1$ is a homomorphism
from $\operatorname{Pic}(B)$ to $I$.  Determinant is its inverse
on the generating line differences, proving (NK.1).  For a line
element, $\gamma_z([M]-1)=1+([M]-1)z$.  Expanding the gamma series
of sums and differences of these generators shows that every term
of gamma degree at least two is a product of two elements of $I$,
and hence is zero.  Finally
$\psi^q([M])=[M]^q=1+q([M]-1)$ proves (NK.2) directly.  This is
an operator calculation, with no semisimplicity assumption.

\emph{3. Injectivity after a normal modification.}
For the moment only assume that $R$ is a normal domain and
$g:Z\to\operatorname{Spec}R$ is proper and birational with $Z$
integral.  A regular function on $Z$ lies in $\operatorname{Frac}R$.
For every height-one prime of $R$, properness extends the generic
point map across the corresponding DVR.  Pulling back the function
shows that it lies in that DVR.  The intersection characterization
of a noetherian normal domain gives $g_*\mathcal O_Z=\mathcal O_R$.
Equivalently, proper coherent finiteness makes $\Gamma(Z,\mathcal
O_Z)$ a finite $R$-subalgebra of its fraction field, and normality
gives equality.  Apply the same argument on affine opens of the base.

Let $g_P:Z\times P\to B$.  On an affine open of $P$, a finite
affine \v{C}ech equalizer for the separated scheme $Z$ computes
global functions.  Tensoring over the field $k$ preserves this
equalizer, so the structure-sheaf equality persists after product
with $P$.  The projection formula for an invertible $M$ gives
\begin{equation}
 g_{P*}\mathcal O_{Z\times P}=\mathcal O_B,
 \qquad g_{P*}g_P^*M=M,
 \qquad\operatorname{Pic}(B)\hookrightarrow
                          \operatorname{Pic}(Z\times P).
                                                        \tag{NK.5}
\end{equation}
If a $K_0(B)$ class pulls back to zero, its rank is zero and its
determinant pulls back to the trivial bundle.  Picard injectivity
and (NK.1) therefore prove
\begin{equation}
                 K_0(B)\hookrightarrow K_0(Z\times P).
                                                        \tag{NK.6}
\end{equation}
No rank-and-determinant assertion about the target is needed.

\emph{4. Lift the whole prescribed $K_1$ class.}
Now impose the characteristic-$p$ regular-modification hypotheses
in the theorem.  The every-degree fixed-product abstract-blow-up
sequence established in FF.1--FF.2 gives
\begin{equation}
\begin{split}
 K_1(Z\times P)[1/p]\oplus K_1(P)[1/p]
 &\longrightarrow K_1(E\times P)[1/p]\\
 &\longrightarrow K_0(B)[1/p]\\
 \longrightarrow K_0(Z\times P)[1/p]\oplus K_0(P)[1/p].
\end{split}                                               \tag{NK.8}
\end{equation}
For clarity, its inputs are the abstract-blow-up pro-exact sequence
\citep[Theorem A and (1.6)]{rootKerzStrunkTamme2018}, nilpotent
invariance after inverting $p$
\citep[Theorem A]{rootGeisserHesselholt2011}, and finite-cover
Zariski descent \citep[Theorem 8.1]{D_ThomasonTrobaugh1990}.
As in FF, nilpotent invariance makes the infinitesimal terms
constant pro-groups, so their ordinary sequence is exact; no inverse
limit or uniform bound on nilpotent exponents is asserted.
By (NK.6) the last map is injective.  Thus the first map is
surjective.  Its $K_1(P)$ contribution is already the restriction
of the actual pullback along $Z\times P\to P$.  Subtracting that
pullback proves (NK.7) on the original fixed product.

\emph{5. Pointwise kernel and a literal endpoint correction.}
The point version of the same sequence gives
\begin{equation}
 \ker\bigl(K_1(Z)_{\mathbf Q}\to K_1(E)_{\mathbf Q}\bigr)
   \subseteq\operatorname{im}\bigl(K_1(R)_{\mathbf Q}
                                      \to K_1(Z)_{\mathbf Q}\bigr),
                                                        \tag{NK.9}
\end{equation}
and, with the residue condition retained, the exact equality
\begin{equation}
\begin{split}
 \ker\bigl(K_1(Z)_{\mathbf Q}\to K_1(E)_{\mathbf Q}\bigr)
 =\operatorname{im}\bigl(&\ker[K_1(R)_{\mathbf Q}\to K_1(k)_{\mathbf Q}]\\
                         &\longrightarrow K_1(Z)_{\mathbf Q}\bigr).
\end{split}                                               \tag{NK.9a}
\end{equation}
Indeed, apply exactness to the pair $(a,0)$ in
$K_1(Z)_{\mathbf Q}\oplus K_1(k)_{\mathbf Q}$; its preimage in
$K_1(R)_{\mathbf Q}$ has residue zero.  The converse follows from
the same diagram.  One must not replace the residue-zero subgroup
by all units.  For a commutative local ring, stable elementary row
reduction gives $K_1(R)=R^*$
\citep[Chapter III, Lemma 1.4 and proof]{frNKWeibelLocalKone}.
On units $\psi^q$ is multiplication by $q$ in additive notation
\citep[Chapter IV, Example 5.4.1]{frNKWeibelAdams}.
Thus the kernel in (NK.9a) is pure weight one.

Assume in addition that $R$ is henselian.  By FF.4, restriction
$K_0(Z)_{\mathbf Q}\to K_0(E)_{\mathbf Q}$ is an isomorphism.
Let $\beta$ be a weight-two class on $E$, and denote its unique
weight-two lift by $\beta_Z$.  Suppose an actual proper family
$\eta\in K_1(E\times P)_{\mathbf Q}$ with points $o,t\in P(k)$
satisfies the literal endpoint equality
$\eta_t-\eta_o=\{u\}\beta$, where $u\in k^*$.
Lift $\eta$ by (NK.7).  Its endpoint difference minus
$\{u\}\beta_Z$ lies in (NK.9a).  Products commute with Adams
operations \citep[Chapter IV, Theorem 5.9 and proof]{frNKWeibelAdams},
so $\{u\}\beta_Z$ has eigenvalue $8$ for $\psi^2$, whereas this
kernel has eigenvalue $2$.  Applying the polynomial
\begin{equation}
                              \frac{\psi^2-2}{6}          \tag{NK.10}
\end{equation}
to the lifted actual family therefore makes its endpoint difference
exactly $\{u\}\beta_Z$.  The natural operations on the schemes
and restrictions are those already used in C.3bt
\citep[p.~258, paragraph preceding Lemma 2.1]{frC3bcLevine1994}.
No projector on the entire $K_1$ group or semisimplicity of that
group is being assumed.
\qed

\emph{Scope of the auxiliary theorem.}
The coefficient ring remains the original local $R$, and the
parameter is the actual fixed proper curve over $k$.
For a new symbolic $u$, the ring $R\otimes_k k(u)$ is generally
neither local nor henselian, so the closed-fibre arguments above
cannot be repeated there.  Henselizing at a chosen point changes
the product.  A class defined on an open of a finite normalization
extends only when its supported localization boundary vanishes;
conjugation carries this boundary to its conjugates and does not
prove their trace zero.  Coherent closure of a class does not
extend a $K_1$ automorphism together with its inverse.  No such
boundary cancellation is asserted here.  Moreover (NK.10) assumes
an exact endpoint on the whole exceptional scheme, not only on
its components.  The component Albanese-kernel quotient of C.3bt
therefore remains unchanged by this auxiliary theorem.
